% Part: methods
% Chapter: proofs
% Section: starting-proofs

\documentclass[../../../include/open-logic-section]{subfiles}

\begin{document}

\olfileid{mth}{prf}{str}

\olsection{सिद्धतेची सुरुवात}

पण सुरुवात करायची तरी कुठून?

तुम्हाला काहीतरी सिद्ध करायला दिले आहे. म्हणून सिद्धतेत शेवटी त्याचाच उल्लेख आला पाहिजे. अर्थात, सुरुवातीलाच आपण ते सिद्ध करणार आहोत असे \emph{जाहीर} करता येते; पण ते गृहीतक म्हणून वापरायचे नाही. कोऱ्या कागदाच्या तळाशी काय सिद्ध करायचे आहे ते लिहा. त्यामुळे तुमचे ध्येय नजरेआड होणार नाही.

यानंतर, वापरता येण्याजोगी काही गृहीतके असतील. तुम्ही करत असलेल्या सिद्धतेच्या \emph{प्रकारा}विषयी पुढील विभागांत बोलताना हे अधिक स्पष्ट होईल. ही गृहीतके पानाच्या वरच्या बाजूला लिहा आणि ती गृहीतके आहेत हे स्पष्टपणे दर्शवा. उदाहरणार्थ, $p$ गृहीत धरत असाल, तर “$p$ असे गृहीत धरा” किंवा “$p$ असे समजा” असे लिहा. शेवटी, प्रश्नात काही व्याख्या असतील ज्या तुम्हाला समजणे आवश्यक आहे. एखादी विशिष्ट व्याख्या वापरायला सांगितलेली असेल; किंवा गृहीतकांत आणि ज्या निष्कर्षापर्यंत पोहोचायचे आहे त्यात विविध व्याख्या आलेल्या असतील. \emph{या व्याख्या लिहून काढा आणि त्यांचा अर्थ तुम्हाला समजतो आहे याची खात्री करा.}

सिद्धतेची मांडणी कशी करायची हे प्रश्नाच्या रूपावरही अवलंबून असते. वाक्याच्या प्रकारानुसार सिद्धतेची मांडणी कशी करायची याचा तपशील पुढील विभागांत दिला आहे.

\end{document}
